%-----------------------------------------------------
% Happy Writing !
%-----------------------------------------------------

% \documentclass[letterpaper, 10 pt, conference]{IEEEtran}
\documentclass[letterpaper, 10 pt, conference, onecolumn]{IEEEtran}
\IEEEoverridecommandlockouts % This command is only needed if you want to use the \thanks command

% Maths
\usepackage{amssymb}
\usepackage{tensor}
\usepackage{cancel}
\usepackage{algorithm}
\usepackage{mathtools}

% Tables
\usepackage{tabularx}
\usepackage{booktabs} % well-spaced, horizontal lines
\usepackage{makecell}
\usepackage{multirow}

% % Figures
\usepackage{float}
\usepackage{graphicx}
\usepackage[caption=false,font=normalsize,labelfont=sf,textfont=sf]{subfig}

% Latex Macro & Utils
\usepackage{xparse} % macros with optional args ande default values
\usepackage{xcolor}
\usepackage{framed} % callout-block background
\usepackage{hyperref}
\usepackage[acronym]{glossaries}
\usepackage{pifont} % get good-looking marks
\usepackage{lipsum} % random text


% theorem environments
\newtheorem{theorem}{Theorem}[section]
\newtheorem{lemma}[theorem]{Lemma}
\newtheorem{corollary}[theorem]{Corollary}
\newtheorem{proposition}[theorem]{Proposition}
\newtheorem{conjecture}[theorem]{Conjecture}
\newtheorem{assumption}[theorem]{Assumption}
\newtheorem{condition}[theorem]{Condition}
\newtheorem{definition}[theorem]{Definition}
\newtheorem{example}[theorem]{Example}
\newtheorem{remark}[theorem]{Remark}

% Shortcuts
\newcommand{\tabref}[1]{Table~\ref{#1}}
\newcommand{\figref}[1]{Fig.~\ref{#1}}
\newcommand{\cmark}{\ding{51}}%
\newcommand{\xmark}{\ding{55}}%

\newcommand{\eg}{e.g., }
\newcommand{\ie}{i.e., }
\newcommand{\etal}{\textit{et al.}} % All the packages and some of the general newcommands
%% This file contains the explicit notation definitions
%% DO NOT import packages within this file.

%---------------------------------------------------------
% Spatial Transformation
%---------------------------------------------------------

% Frame
\providecommand{\frameX}[1]{\{#1\}}

% Transformation 
% Usege: \tf[1][2]{3}[4][5]. If one optional argument appears on either side of var, it will be considered as the first argument on that side.
% #1 (optional): where the quantity is expressed w.r.t or the coordinates
% #2 (optional): where quantity is the measured w.r.t or reference.
% #3 (mendatory): the symbol of the transformation, typically R, t, X
% #4 (optional): the tip or index.
% #5 (optional): the operator or superscript - usually \top, or \vee, or \times, etc.

\NewDocumentCommand{\tf}{O{} O{} m o O{}}  
  {
    \tensor*[^{#1}_{#2}]{\mathbf{#3}}{_{#4}^{#5}}
  }

%---------------------------------------------------------
% Sets
%---------------------------------------------------------

\providecommand{\C}{\mathbb{C}} % complex Number
\providecommand{\N}{\mathbb{N}} % Natural
\providecommand{\R}{\mathbb{R}} % Real
\providecommand{\Z}{\mathbb{Z}} % Integer

% Lie groups
\providecommand{\SO}{\mathbf{SO}}
\providecommand{\SL}{\mathbf{SL}}
\providecommand{\GL}{\mathbf{GL}}
\providecommand{\SE}{\mathbf{SE}}
\providecommand{\SIM}{\mathbf{SIM}}
\providecommand{\grpG}{\mathbf{G}}
\providecommand{\grpH}{\mathbf{H}}
\providecommand{\grpK}{\mathbf{K}}
\providecommand{\grpN}{\mathbf{N}}
\providecommand{\grpR}{\mathbf{R}}
\providecommand{\grpS}{\mathbf{S}}
\providecommand{\grpZ}{\mathbf{Z}}

% Lie algebras
\providecommand{\gothgl}{\mathfrak{gl}}
\providecommand{\gothso}{\mathfrak{so}}
\providecommand{\gothse}{\mathfrak{se}}
\providecommand{\gothsl}{\mathfrak{sl}}
\providecommand{\gothsim}{\mathfrak{sim}}
\providecommand{\gothsot}{\mathfrak{sot}}
\providecommand{\gothmr}{\mathfrak{mr}}
\providecommand{\gothe}{\mathfrak{e}}
\providecommand{\gothg}{\mathfrak{g}}
\providecommand{\gothh}{\mathfrak{h}}
\providecommand{\gothk}{\mathfrak{k}}
\providecommand{\gothl}{\mathfrak{l}}
\providecommand{\gothn}{\mathfrak{n}}
\providecommand{\gothr}{\mathfrak{r}}
\providecommand{\goths}{\mathfrak{s}}
\providecommand{\gothu}{\mathfrak{u}}
\providecommand{\gothv}{\mathfrak{v}}
\providecommand{\gothz}{\mathfrak{z}}
\providecommand{\gothB}{\mathfrak{B}}
\providecommand{\gothK}{\mathfrak{K}}
\providecommand{\gothL}{\mathfrak{L}}
\providecommand{\gothR}{\mathfrak{R}}
\providecommand{\gothU}{\mathfrak{U}}
\providecommand{\gothV}{\mathfrak{V}}
\providecommand{\gothW}{\mathfrak{W}}
\providecommand{\gothX}{\mathfrak{X}} % as in X(M)
\providecommand{\gothZ}{\mathfrak{Z}}

%~~~~~~~~~~~~~~~~~~~~~~~~~~~~~~~
% shortcuts
\providecommand{\so}{\mathfrak{so}}
\providecommand{\se}{\mathfrak{se}}
\providecommand{\sot}{\mathfrak{sot}}
%~~~~~~~~~~~~~~~~~~~~~~~~~~~~~~~

% vector spaces
\providecommand{\vecG}{\mathbb{G}} % tangent space at X_0
\providecommand{\vecK}{\mathbb{K}} % kernel of linear operator.
\providecommand{\vecO}{\mathbb{O}}
\providecommand{\vecU}{\mathbb{U}}
\providecommand{\vecV}{\mathbb{V}}
\providecommand{\vecW}{\mathbb{W}}
\providecommand{\vecY}{\mathbb{Y}}
\providecommand{\bfV}{\mathbf{V}}
\providecommand{\calV}{\mathcal{V}}

% Euclidean space
\providecommand{\eucE}{\mathbb{E}}

% Matrix spaces
\providecommand{\Sym}{\mathbb{S}} % symmetric matrix $\Sym(n)$
\providecommand{\PD}{\mathbb{S}_+} % positive definite matrices.

% Frame bundle
\providecommand{\FSO}{\calF_{T \SO}}
\providecommand{\FB}{\calF} % framebundle.  But also extended function.

%---------------------------------------------------------
% functions and mapping
%---------------------------------------------------------

% operators
\DeclareMathOperator{\inv}{inv}
\DeclareMathOperator{\tr}{tr}
\DeclareMathOperator{\vex}{vrp}
\DeclareMathOperator{\grad}{grad}
\DeclareMathOperator{\spn}{span}
\DeclareMathOperator{\diag}{diag}
\DeclareMathOperator{\stab}{stab}
\DeclareMathOperator{\Ad}{Ad}
\DeclareMathOperator{\ad}{ad}
\DeclareMathOperator{\kernel}{ker}
\DeclareMathOperator{\image}{im}
\DeclareMathOperator*{\argmin}{argmin}
\DeclareMathOperator*{\argmax}{argmax}

% maps
\providecommand{\id}{\mathrm{id}} % identity map
\providecommand{\pr}{\mathbb{P}} % projection
\providecommand{\prse}{\mathbb{P}_{\se}} % projection
\providecommand{\mm}{\mathbb{m}} % monomorphism

% cost
\providecommand{\cost}{\ell} %% local costs
\providecommand{\Lyap}{\mathcal{L}} %% aggregate cost

%---------------------------------------------------------
% Differential notation
%---------------------------------------------------------

% differentials
\providecommand{\tT}{\mathrm{T}} % tangent bundles
\providecommand{\td}{\mathrm{d}}
\providecommand{\tD}{\mathrm{D}}
\providecommand{\tL}{\mathrm{L}}
\providecommand{\pd}{\partial}
\providecommand{\ddt}{\frac{\td}{\td t}}
\providecommand{\Ddt}{\frac{\tD}{\td t}}
\providecommand{\ddtau}{\frac{\td}{\td \tau}}
\providecommand{\pdt}{\frac{\partial}{\partial t}}
\providecommand{\pdtau}{\frac{\partial}{\partial \tau}}
\providecommand{\dt}{\td t}
\providecommand{\dtau}{\td \tau}
\providecommand{\Hess}{\mathrm{Hess}}


%---------------------------------------------------------
% Variables
%---------------------------------------------------------

% group elements
\providecommand{\eye}[1]{\mathbf{I}_{#1}} % identity of a matrix group.

% vector space elements
\providecommand{\eb}{\mathbf{e}} % \providecommand{\eb}{\vec{e}}





 % For general notations defs

%-----------------------------------------------------
% Front Matter
%-----------------------------------------------------

\begin{document}

\title{\LARGE \bf The Best Title (with A Catchy Acronym)}
\author{Author A$^{1}$ and Author B$^{2}$
    \thanks{*Author contributions}
    \thanks{$^{1}$AuthorA affiliations and contact info
            {\tt\small authorA@mail.com}
    }
    \thanks{$^{2}$AuthorB affiliations and contact info
            {\tt\small authorB@mail.com}
    }
    \thanks{This project is funded by ...}
}
\maketitle

\begin{abstract}
    \lipsum[1-3]
\end{abstract}

%-----------------------------------------------------
% Main Matter
%-----------------------------------------------------

\section{Introduction}

\lipsum[1-3]


%-----------------------------------------------------

\section{Related Work}
\lipsum[1-3]

\begin{figure*}[!t]
    \centering
    \subfloat[]{\includegraphics[width=2.5in]{asset/placeholder_image.jpeg}%
        \label{fig_first_case}}
    \hfil
    \subfloat[]{\includegraphics[width=2.5in]{asset/placeholder_image.jpeg}%
        \label{fig_second_case}}
    \caption{Dae. Ad quatur autat ut porepel itemoles dolor autem fuga. Bus quia con nessunti as remo di quatus non perum que nimus. (a) Case I. (b) Case II.}
    \label{fig_sim}
\end{figure*}

\lipsum[1-3]


%-----------------------------------------------------

\section{Methodology}

Antonova \etal~\cite{antonova_2021_SequentialTopologicalRepresentations} proposed bla.

\begin{table*}
    \caption[Comparison of runtime and multi-instance capability between popular pose-estimation algorithms.]{Comparison of runtime and multi-instance capability between popular pose-estimation algorithms.}
    \centering
    \begin{tabular}{c c c c c c c}
        \toprule
                       & \makecell{PVNet$^*$ \\ABC} & \makecell{PoseCNN +\\ICP~} & \makecell{DOPD\\ ABC} & \makecell{DenseFusion\\} & \makecell{SSD6D+\\ICP~} & \makecell{DOPE} \\
        \midrule
        Modality       & RGB                 & RGB+D  & RGB    & RGBD   & RGB+D  & RGB \\
        \midrule
        Multi-Instance & \xmark              & \xmark & \xmark & \xmark & \cmark & \cmark \\
        \midrule
        fps            & 25                  & 0.1    & 33     & 16     & 10     & 10 \\
        \bottomrule
        \multicolumn{7}{l}{\rule{0pt}{3.5ex} \textit{\makecell[l]{$^*$The reported runtime of PVNet is the based on the conference version, which does not incorporate \\ the instance-level segmentation pre-processing pipeline.}}}
    \end{tabular}
    \label{tab:pose_compare}
\end{table*}

\lipsum[1-3]





%-----------------------------------------------------

\section{Experiments}

\begin{table}
    \caption{Per action recognition performance.}
    \centering
    \begin{tabular}{c c c c}
        \toprule
        \textbf{Action}    & \textbf{precision} & \textbf{recall} & \textbf{f1-score} \\
        \midrule
        no assembly action & 0.72               & 0.80            & 0.75 \\
        reach              & 0.40               & 0.13            & 0.14 \\
        flip table top     & 0.79               & 0.80            & 0.79 \\
        flip table         & 0.64               & 0.52            & 0.53 \\
        spin leg           & 0.79               & 0.61            & 0.66 \\
        align leg          & 0.48               & 0.37            & 0.40 \\
        rotate table       & 0.66               & 0.56            & 0.56 \\
        human grasp        & 0.27               & 0.17            & 0.13 \\
        \midrule
        Average            & 0.59               & 0.49            & 0.50 \\
        \bottomrule
    \end{tabular}
    \label{tab:results:per_action_performance}
\end{table}



%-----------------------------------------------------

\section{Conclusion}


%-----------------------------------------------------
% References
%-----------------------------------------------------
\bibliographystyle{util/IEEEtranS}
\bibliography{bib/IEEEabrv, bib/Deformable-Objects}

\end{document}